% COMSYS Seminar/Preseminar Paper Template
% ----------------------------------------
% Based on the original IEEE template package, includes annotations by COMSYS

\documentclass[conference]{IEEEtran}

%\IEEEoverridecommandlockouts
\usepackage{amsmath,amssymb,amsfonts}
\usepackage{algorithmic}
\usepackage[hidelinks]{hyperref}
\usepackage[capitalize,nameinlink]{cleveref}
\usepackage[sort]{cite}
\usepackage[inline]{enumitem}
\usepackage{balance}
\usepackage{graphicx}
\usepackage{textcomp}
\usepackage{booktabs}
\usepackage{microtype}



% Prevent subcaption package from overwriting the IEEE template defaults
\makeatletter
\let\SafeCaption\@makecaption
\makeatother
\usepackage{subcaption}
\makeatletter
\let\@makecaption\SafeCaption
\makeatother

\def\BibTeX{{\rm B\kern-.05em{\sc i\kern-.025em b}\kern-.08em
    T\kern-.1667em\lower.7ex\hbox{E}\kern-.125emX}}

\begin{document}

\title{Secure Multiparty Communication Systems}

\author{\IEEEauthorblockN{Finn Paul Lippok}
\IEEEauthorblockA{\textit{Seminar on Internet Technology} \\
\textit{Chair of Communication and Distributed Systems}\\
\textit{RWTH Aachen University}\\
Aachen, Germany \\
finn.lippok@rwth-aachen.de}
}

\maketitle

\begin{abstract}
Placeholder.
\end{abstract}

\begin{IEEEkeywords}
Secure Multiparty Communication
\end{IEEEkeywords}

\section{Introduction}

\begin{itemize}
    \item Motivation: groups are harder than pairs.
    \item Which techniques and protocols are most relevant in current systems? $\rightarrow$ \Cref{sec:techniques}
    \item How practical are these to deploy? $\rightarrow$ \Cref{sec:practicality}
    \item How to protect against misbehaving peers? $\rightarrow$ \Cref{sec:misbehavior}
\end{itemize}

\section{Background}\label{sec:background}
\subsection{Security Goals}
We consider

\subsubsection{Confidentiality}
Sent group messages should only be decryptable by current members of the group.

\subsubsection{Group Membership Authentication}

\subsubsection{Availability}
A third party should not be able to disrupt message delivery or key agreement.
Group members should not be able to unilaterally split the group or disrupt the functionality of other group members.

\begin{itemize}
    \item Sender Authentication
    \item Forward Secrecy
    \item Backward Secrecy
    \item Post-Compromise Security
    \item \textit{Note:} Clear up definition of Member join / leaving FS/BS vs messaging FS/PCS \cite{prantlSurveySecureGroup2022}.
    \item Consistency
    \item Out of Scope(?): Group membership privacy, metadata.
\end{itemize}

\subsection{Threat Model}
\begin{itemize}
    \item Network adversary: passive, active
    \item Malicious Server
    \item Member: state compromise, malicious insider.
\end{itemize}

\subsection{Two-Party End-to-End Encryption}
The most widely deployed protocol for two-party messaging is the Signal Protocol~\cite{balbasWhatsUppSenderKeys2023}, which uses PQXDH~\cite{kretPQXDHKeyAgreement} for initial key agreement and the Double Ratchet Algorithm~\cite{perrinDoubleRatchetAlgorithm}.
In the following, we outline how the protocol works and how it provides FS and PCS (called break-in recovery in the specification~\cite{perrinDoubleRatchetAlgorithm}).
For simplicity, we omit key authentication and the recently added post-quantum ratchet~\cite{perrinDoubleRatchetAlgorithm}.
\subsubsection{Setup}
We assume Alice wants to initiate a conversation with Bob while Bob is offline.
Bob has previously published a prekey bundle to a server, which contains Bob's public key. Alice fetches it and uses the PQXDH key agreement protocol~\cite{kretPQXDHKeyAgreement} to derive a root key.
In the first message, Alice includes the values Bob needs to derive the same root key.
\subsubsection{Symmetric-key ratchet}
Alice and Bob use a KDF to further derive sending and receiving chain keys. For each message they send or receive, they derive a one-time message key, used to encrypt or decrypt the message, and a new sending or receiving chain key, respectively.
Because the KDF is one-way and all old keys are deleted after use, an attacker that compromises the state of either or both parties cannot compute previous chain keys and thus cannot derive the keys used to encrypt old messages.
This provides FS. PCS, on the other hand, is not provided because an attacker can advance the chain exactly as the legitimate parties do and thus derive all future keys~\cite{perrinDoubleRatchetAlgorithm}.
\subsubsection{Diffie--Hellman ratchet}
To provide PCS, new randomness needs to be mixed into the keys. This is achieved by the DH ratchet.
Whenever the sending direction changes, the new sender includes a new ratchet public key in their first message.
From that and the current root key, a new shared root key and chain keys can be derived for the following messages~\cite{perrinDoubleRatchetAlgorithm}.

\subsection{Client-Side Fanout}
The naive approach to scale the existing Signal Protocol to group chats would be client-side fanout.
By creating pairwise channels between all group members, a de facto group chat is created.
For a group of size $n$, each individual is required to create $n-1$ channels and every sent message needs to be encrypted $n-1$ times~\cite{cohn-gordonEndstoEndsEncryptionAsynchronous2018}.
Because every message is encrypted independently, malicious group members could send completely different messages to different members of the group without being caught~\cite{goldbergMultipartyOfftherecordMessaging2009}.
Detecting such misbehaving peers would require additions to the protocol, like making all parties share the hash of their transcripts with each other to catch any discrepancies, as described by Goldberg et al.~\cite{goldbergMultipartyOfftherecordMessaging2009}.
On the other hand, client-side fanout inherits FS and PCS for the same reason~\cite{cohn-gordonEndstoEndsEncryptionAsynchronous2018,roslerMoreLessEndtoEnd2018}.


\section{Techniques in Current Systems}\label{sec:techniques}
In this section, we will discuss protocols used in current systems, their current shortcomings, and proposed improved variants.
We will first discuss the currently most widely adopted approach, Sender Keys, and then go into more detail on the current research direction of Continuous Group-Key Agreement (CGKA), with a focus on tree-based approaches, while also briefly discussing other approaches.
\subsection{Sender Keys}
Sender Keys are currently adopted by many group messaging platforms, most notably WhatsApp and Signal~\cite{balbasWhatsUppSenderKeys2023}.
We will first discuss their functionality as formally described by Balb\'as et al.~\cite{balbasWhatsUppSenderKeys2023} and then \ldots
Like client-side fanout, Sender Keys rely on pairwise channels between group members, but unlike client-side fanout, they do not exchange messages over these channels directly~\cite{balbasWhatsUppSenderKeys2023}.
Instead, a group member $i$ who wants to send a message must have previously sent their sender key to every other group member over such a two-party channel.
Such a sender key consists of a symmetric chain key $\mathsf{ck}_i$ used to derive message keys and a signing public key $\mathsf{pk}_i$, while the corresponding signing $\mathsf{sk}_i$ remains secret to member $i$.
When sending a message, it is encrypted with a message key derived from the current $\mathsf{ck}_i$ and signed with $\mathsf{sk}_i$.
All group members receive the same ciphertext, verify its signature using $\mathsf{pk}_i$ and decrypt it by also deriving the same current message key using their copy of $\mathsf{ck}_i$.
After sending and receiving, respectively, each party performs one ratchet step on $\mathsf{ck}_i$ using an KDF.

Because Sender Keys only use a symmetric chain $\mathsf{ck}_i$ and no second ratchet, PCS can only be achieved by redistributing new sender keys and thus new chain keys.
Because this is implementation-dependent and often only done after a group member leaves, PCS is only provided over long periods of time~\cite{balbasWhatsUppSenderKeys2023}.
If an attacker also compromises the state of the pairwise channels and the new keys are distributed before these channels have fully healed, the newly distributed sender keys could also leak~\cite{balbasWhatsUppSenderKeys2023}.
\begin{itemize}
    \item Similar: Megolm (Matrix)~\cite{balbasWhatsUppSenderKeys2023}
    \item complexity
\end{itemize}

\subsection{Continuous Group-Key Agreement}
Well-designed CGKAs improve upon the weakness of Sender Keys by no longer requiring keys to be shared \ldots with every group member \ldots
\subsubsection{Definitions}
Epochs, commits, Operations
\subsubsection{Tree Based}
\begin{itemize}
    \item Predecessors (Not Strictly CGKA)
    \begin{itemize}
        \item LKH: Server holds all keys
        \item TGDH: not asynchronous
    \end{itemize}
    \item ART \cite{cohn-gordonEndstoEndsEncryptionAsynchronous2018}
    \item TreeKEM/MLS \cite{barnesMessagingLayerSecurity2023} (used by RCS, Discord VC)
    \begin{itemize}
        \item Improved variants:
        \begin{itemize}
            \item RTreeKEM: updates private keys (UPKE) \cite{alwenSecurityAnalysisImprovements2020}
            \item Tainted TreeKEM: avoids blanks \cite{kleinKeepDirtTainted2021}
        \end{itemize}
    \end{itemize}
\end{itemize}
\subsubsection{Other Approaches}
\begin{itemize}
    \item Chained CmPKE \cite{hashimotoConcreteTreatmentEfficient2021}
    \item DCGKA: decentralized (no ordering server) \cite{weidnerKeyAgreementDecentralized2021}
\end{itemize}

\begin{table*}[htbp]
    \centering
    \caption{Comparison}
    \label{tab:comparison}
    \begin{tabular}{lccccccccc}
        \toprule
        Scheme & FS & PCS & Asynchronous & Ordering requirement & Add & Remove & Send & Receive & Update \\
        \midrule
        Fanout & & & & & & & & & \\
        Sender Keys & & & & & & & & & \\
        ART & & & & & & & & & \\
        TreeKEM & & & & & & & & & \\
        RTreeKEM & & & & & & & & & \\
        Tainted TreeKEM & & & & & & & & & \\
        CmPKE & & & & & & & & & \\
        DCGKA & & & & & & & & & \\
        \bottomrule
    \end{tabular}
\end{table*}

\section{Practicality of Deployment}\label{sec:practicality}
\subsection{Complexity}
\begin{itemize}
    \item see \Cref{tab:comparison}
    \item Best vs.\ worst case
    \item Bandwidth asymmetry \cite{hashimotoConcreteTreatmentEfficient2021}
\end{itemize}

\subsection{Other Challenges}
\begin{itemize}
    \item Post-Quantum
    \item Forks and state loss
    \item Interoperability (EU DMA)
    \item Dependence on the ordering server.
\end{itemize}


\section{Protecting Against Misbehaving Peers}\label{sec:misbehavior}
\begin{itemize}
    \item What insiders cannot be prevented from doing
    \item Inconsistent commits $\rightarrow$ split group. (prevent with Zero-Knowledge Proofs \cite{alwenContinuousGroupKey2020})
    \item Insider DoS: abort vs.\ guaranteed output delivery \cite{dalskovFastFullySecure2022}
    \item Cheater identification
\end{itemize}

\section{Conclusion}
% Balance columns on last page of the paper
\balance

% Modified so that bibliography resides in external file.
\bibliographystyle{IEEEtran}
\bibliography{../references}

\end{document}
